\section{Concluding Remarks}\label{sec:conclusion}

This study examines the identification of Smart Money in Polymarket using historical forecasting performance and wallet trading behavior. By combining statistical inference, Empirical Bayes estimation, predictive modeling, and behavioral analysis, it investigates whether forecasting ability persists over time and whether individual wallet activity provides information beyond market probabilities. The results reveal statistically significant persistence in forecasting performance, particularly among experienced wallets. Empirical Bayes estimation further improves the selection of wallets with stronger subsequent forecasting performance. Historical trading behavior also exhibits potential incremental predictive value, demonstrating the usefulness of incorporating behavioral information into probability forecasting.

A particularly interesting finding concerns changes in betting size. Although wallets tend to increase their stakes on outcomes already favored by the market, the additional forecasting advantage associated with larger bets is most pronounced when market probabilities are close to 50\%. This relationship highlights the potential importance of market uncertainty in understanding the informational value of wallet behavior. These findings suggest that identifying Smart Money involves more than evaluating historical forecasting records. The circumstances under which wallets change their trading behavior may also reveal valuable predictive information. The analysis provides a foundation for further research on information leadership, behavioral signals, and trading intelligence in prediction markets.
